\begin{exercise}{41}
Show that the coherent nerve functor $\mathrm{N}\colon\mathrm{Cat}_{\Delta}\to\mathrm{sSet}$ commutes with
coproducts. Show that $\mathfrak C$ is not right-adjoint to $\mathrm{N}$. Does $\mathrm{N}$ have a right adjoint at all? Hint:
Does the ordinary nerve functor $\mathrm{N}\colon\mathrm{Cat}\to\mathrm{sSet}$ have a right adjoint? And how are these two
questions related?

\begin{proof}
The coproduct $\coprod_i\mathcal C_i$ has the disjoint union
of the object sets, the original hom-simplicial sets within
each summand, and empty hom-simplicial sets between summands.
Every simplicial functor
$\mathfrak C[\Delta^n]\to\coprod_i\mathcal C_i$ lands in a
unique summand: its objects $0,\ldots,n$ are joined by
successive hom-simplicial sets $\Delta^0$, so consecutive
images must belong to the same summand.
For $n=0$, its sole object determines the summand.
By Definition~1.2.63 this gives bijections
\[
 \mathrm{N}\Bigl(\coprod_i\mathcal C_i\Bigr)_n
 \cong\coprod_i\mathrm{N}(\mathcal C_i)_n.
\]
They commute with precomposition by every map in $\Delta$,
and hence prove the asserted coproduct isomorphism.

Write $\mathrm{N}_{\mathrm{ord}}$ for the ordinary nerve.
In $\mathrm{Cat}$ the coequalizer of the two object inclusions
$[0]\rightrightarrows[1]$ is $B\mathbb N$, the one-object
category whose endomorphism monoid is $(\mathbb N,+)$.
Indeed, functors from $B\mathbb N$ to a category $\mathcal A$
are exactly an object of $\mathcal A$ and an endomorphism of it,
which is also the data of a functor $[1]\to\mathcal A$
agreeing on the two objects.
The coequalizer $Q$ of the two vertex inclusions
$\Delta^0\rightrightarrows\Delta^1$ is computed degreewise.
Its set $Q_2$ has three elements: the four monotone maps
$[2]\to[1]$, with the two constant maps identified.
In contrast,
\[
 \mathrm{N}_{\mathrm{ord}}(B\mathbb N)_2
 \cong\mathbb N\times\mathbb N
\]
is infinite. Thus $\mathrm{N}_{\mathrm{ord}}$ does not
preserve this coequalizer and has no right adjoint.

Finally, $c\dashv u$ by Corollary~1.2.51, and
$\mathrm{N}\circ c\cong\mathrm{N}_{\mathrm{ord}}$ by
Lemma~1.2.64. If $\mathrm{N}$ had a right adjoint $R$, then
$uR$ would be right-adjoint to $\mathrm{N}_{\mathrm{ord}}$,
a contradiction. In particular, $\mathfrak C$ is not
right-adjoint to $\mathrm{N}$.
\end{proof}
\end{exercise}
